\section{Scalars from gates with no participants}
\label{sec:peps_empty_gate_scalars}

\begin{definition}[Empty-participant scalar]
    \label{def:peps_empty_gate_scalar}
    \lean{TNLean.PEPS.PairEffect.SourceCircuit.emptyGateScalar}
    \leanok
    For each gate with no participants, sum its original monomial
    coefficients. Let $s$ be the product of these sums over the
    chronological circuit. All other operations contribute the factor $1$.
    An empty monomial family contributes $0$.
\end{definition}

\begin{theorem}[The scalar part of the chronological residual]
    \label{thm:peps_empty_gate_scalar_residual}
    \lean{TNLean.PEPS.PairEffect.SourceCircuit.norm_emptyGateScalar_le_one,
      TNLean.PEPS.PairEffect.SourceCircuit.eval_heq_emptyGateScalar_smul_id_of_isEmpty,
      TNLean.PEPS.PairEffect.SourceCircuit.restrict_partialResidual_all_eval_heq_emptyGateScalar}
    \leanok
    \uses{def:peps_empty_gate_scalar}
    The scalar satisfies $|s|\leq 1$. When all original parties are
    selected, the restriction of the chronological residual to the
    remaining scalar tensor factor is multiplication by $s$.
    In particular, a circuit on an empty party set acts on the empty
    tensor product as multiplication by $s$.
\end{theorem}

\begin{proof}\leanok
    Induct on the chronological circuit. Composition multiplies the
    scalar factors. Operations at selected parties and their source
    registers do not alter the remaining scalar factor. A gate with
    no participants contributes exactly its coefficient sum. Its
    contractivity bounds the modulus of this sum by $1$, and the
    product remains contractive.
\end{proof}

\begin{theorem}[Absorption into an existing party tensor]
    \label{thm:peps_network_scalar_absorption}
    \lean{TNLean.PEPS.DependentBondNetwork.network_mul_at_vertex}
    \leanok
    Let $A$ be a finite tensor network, let $p$ be one of its vertices,
    and let $c\in\mathbb C$. Multiplying the local tensor at $p$ by
    $c$ multiplies the contraction by $c$. The bonds and their dimensions
    are unchanged. This applies also when $c=0$.
\end{theorem}

\begin{proof}\leanok
    At each virtual configuration, the vertex product acquires precisely
    the factor $c$. Pull this factor out of the finite sum.
\end{proof}

These statements justify retaining the scalar from empty-participant gates
until the final absorption step in
\cite[proof of Theorem~5.2]{OpenAI2026PolynomialPEPS},
\texttt{04-compression.tex}, lines 565--588.
The absorption statement specifies an existing vertex; it makes no claim
that an arbitrary scalar is the contraction of a network with no vertices.
